
\documentclass{article}
\usepackage{colortbl}
\usepackage{makecell}
\usepackage{multirow}
\usepackage{supertabular}

\begin{document}

\newcounter{utterance}

\twocolumn

{ \footnotesize  \setcounter{utterance}{1}
\setlength{\tabcolsep}{0pt}
\begin{supertabular}{c@{$\;$}|p{.15\linewidth}@{}p{.15\linewidth}p{.15\linewidth}p{.15\linewidth}p{.15\linewidth}p{.15\linewidth}}

    \# & $\;$A & \multicolumn{4}{c}{Game Master} & $\;\:$B\\
    \hline 

    \theutterance \stepcounter{utterance}  

    & & \multicolumn{4}{p{0.6\linewidth}}{\cellcolor[rgb]{0.9,0.9,0.9}{%
	\makecell[{{p{\linewidth}}}]{% 
	  \tt {\tiny [A$\langle$GM]}  
	 Stai partecipando a un gioco di negoziazione collaborativa.\\ \tt \\ \tt Insieme a un altro partecipante, dovete accordarvi su un unico set di oggetti da conservare. Ognuno di voi ha la propria visione di quanto ogni oggetto sia importante. Non sai come l'altro partecipante valuta gli oggetti. Inoltre, ti viene fornito lo sforzo che ogni oggetto richiede.\\ \tt Potete accordarvi su un set solo se lo sforzo totale degli oggetti selezionati non supera un limite condiviso:\\ \tt \\ \tt LIMITE = 7868\\ \tt \\ \tt Ecco i valori di sforzo dei singoli oggetti:\\ \tt \\ \tt Sforzo oggetto = {"B21": 552, "B00": 560, "A08": 76, "A27": 950, "B94": 248, "B53": 871, "A23": 391, "C65": 143, "A89": 293, "B67": 208, "B45": 678, "C34": 737, "C12": 407, "A05": 366, "B60": 768, "C92": 859, "B69": 183, "A84": 231, "C06": 305, "B12": 727, "A16": 148, "B85": 357, "C35": 504, "B10": 549, "C15": 299, "B09": 91, "B16": 527, "C81": 848, "A58": 307, "C24": 214, "A30": 723, "A66": 475, "B75": 23, "B46": 298, "C70": 820}\\ \tt \\ \tt Ecco la tua visione personale sull'importanza di ogni oggetto:\\ \tt \\ \tt Valori di importanza dell'oggetto = {"B21": 138, "B00": 583, "A08": 868, "A27": 822, "B94": 783, "B53": 65, "A23": 262, "C65": 121, "A89": 508, "B67": 780, "B45": 461, "C34": 484, "C12": 668, "A05": 389, "B60": 808, "C92": 215, "B69": 97, "A84": 500, "C06": 30, "B12": 915, "A16": 856, "B85": 400, "C35": 444, "B10": 623, "C15": 781, "B09": 786, "B16": 3, "C81": 713, "A58": 457, "C24": 273, "A30": 739, "A66": 822, "B75": 235, "B46": 606, "C70": 105}\\ \tt \\ \tt Obiettivo:\\ \tt \\ \tt Il tuo obiettivo è negoziare un set condiviso di oggetti che ti avvantaggi il più possibile (cioè, massimizzi l'importanza totale per TE), rimanendo entro il LIMITEE. Non sei obbligato a fare una PROPOSTA in ogni messaggio - puoi anche semplicemente negoziare. Tutte le tattiche sono permesse!\\ \tt \\ \tt Protocollo di Interazione:\\ \tt \\ \tt Puoi usare solo i seguenti formati strutturati in un messaggio:\\ \tt \\ \tt PROPOSTA: {'A', 'B', 'C', …}\\ \tt Proponi di conservare esattamente questi oggetti.\\ \tt RIFIUTO: {'A', 'B', 'C', …}\\ \tt Rifiuta esplicitamente la proposta dell'avversario.\\ \tt ARGOMENTAZIONE: {'...'}\\ \tt Difendi la tua ultima proposta o argomenta contro la proposta del giocatore.\\ \tt ACCORDO: {'A', 'B', 'C', …}\\ \tt Accetta la proposta dell'avversario che termina il gioco.\\ \tt RAGIONAMENTO STRATEGICO: {'...'}\\ \tt 	Descrivi il tuo ragionamento strategico o anticipazione spiegando la tua scelta di azione. Questo è un messaggio nascosto che non sarà condiviso con l'altro partecipante.\\ \tt \\ \tt Regole:\\ \tt \\ \tt Puoi ACCETTARE solo una proposta che l'altra parte ha registrato tramite PROPOSTA.\\ \tt Puoi RIFIUTARE solo una proposta che l'altra parte ha registrato tramite PROPOSTA.\\ \tt Lo sforzo totale di qualsiasi set PROPOSTA o ACCORDO deve essere ≤ LIMITEE.\\ \tt NON rivelare i tuoi punteggi di importanza nascosti.\\ \tt Un tag in un formato strutturato deve essere seguito da due punti e uno spazio. L'argomento deve essere un set python contenente 0 o più stringhe.\\ \tt Quindi, deve essere della forma TAG: {...}\\ \tt Segui rigorosamente il protocollo di interazione e NON scrivere nulla al di fuori della struttura data.\\ \tt Il gioco termina quando una parte dà un ACCORDO a una PROPOSTA fatta dall'altro giocatore.\\ \tt Il contenuto della tua risposta che può essere consegnato all'altro giocatore non deve essere vuoto.\\ \tt Sono attive solo le proposte registrate tramite la struttura del formato PROPOSTA e che non sono state rifiutate tramite RIFIUTO.\\ \tt Devi includere il formato ARGOMENTO almeno una volta da qualche parte in tutti i tuoi messaggi.\\ \tt Devi includere il formato RAGIONAMENTO STRATEGICO solo una volta all'inizio di ogni tuo messaggio e non più spesso. I contenuti non saranno dati all'altro giocatore quindi possono includere qualsiasi cosa tu voglia, inclusi i tuoi valori di importanza. Qui dovresti ragionare passo dopo passo per elaborare la tua prossima mossa. 
	  } 
	   } 
	   } 
	 & \\ 
 

    \theutterance \stepcounter{utterance}  

    & \multicolumn{4}{p{0.6\linewidth}}{\cellcolor[rgb]{0.8,1,0.9}{%
	\makecell[{{p{\linewidth}}}]{% 
	  \tt {\tiny [A$\rangle$GM]}  
	 RAGIONAMENTO STRATEGICO: {'Obiettivo: proporre un insieme ad alto valore per me rispettando il limite 7868, lasciando margine per eventuali concessioni. Calcolo rapporto valore/sforzo per priorità. Seleziono greedy per valore totale elevato e includo molti oggetti piccoli con alto valore per facilitare scambi. Scelgo: B12(915/727), A08(868/76), A16(856/148), B09(786/91), B60(808/768), C15(781/299), B94(783/248), B67(780/208), A30(739/723), C81(713/848 forse pesante ma alto valore), B10(623/549), B46(606/298), B00(583/560), A66(822/475), A27(822/950 forse pesante), C12(668/407), A89(508/293), A58(457/307), B45(461/678 forse basso), C34(484/737 basso), A05(389/366), B85(400/357), C35(444/504), C24(273/214), C65(121/143), ecc. Calcolo una combinazione entro 7868. Provo a costruire set: Parto con top efficienza: A08 76, A16 148, B09 91, B94 248, B67 208, C15 299, A66 475, A89 293, B46 298, B12 727, B10 549, B00 560, A30 723, B60 768, C12 407, A58 307, A05 366, B85 357, C24 214. Sommo sforzi: 76+148+91+248+208+299+475+293+298+727+549+560+723+768+407+307+366+357+214 = 6714. Valore somma: 868+856+786+783+780+781+822+508+606+915+623+583+739+808+668+457+389+400+273 = 11763. Spazio residuo 1154. Considero aggiunte: C81(848) valore 713 tot 7562; resta 306. Con 306 posso aggiungere C06(305, val30) o C65(143,121) o B69(183,97) o B75(23,235) etc. Provo C06 305 -> sforzo 7867 (ok) valore +30=11793. Ma potrei meglio: usare 306 per A84(231,500) e B75(23,235) totale 254, resta 52; molto meglio. Provo sostituire C06 con A84+B75: nuovo sforzo 7562+231+23=7816; resta 52. Con 52 nessun altro item (minimi: A08 76, B75 23 già preso). Valore +500+235=12498. Ottimo. Controllo set totale: Sforzo 7816 <= 7868. Valore totale 11763 (base senza C81/A84/B75?) Ricomputo: Base 6714 valore 11763; aggiungo C81 val713 sforzo +848 => sforzo 7562 valore 12476; aggiungo A84 val500 sforzo +231 => 7793 valore 12976; aggiungo B75 val235 sforzo +23 => 7816 valore 13211. Questo è molto alto. Verifico che non ho incluso elementi a valore quasi nullo come B16(3). Controllo se sostituire qualcosa per ancora valore: abbiamo spazio 52, nessun aggiunta utile. Ottimo. Set proposto: {A08,A16,B09,B94,B67,C15,A66,A89,B46,B12,B10,B00,A30,B60,C12,A58,A05,B85,C24,C81,A84,B75}. Pronto a difendere con argomentazione su equilibrio sforzo e copertura rischi. Pronto a cedere su alcuni pezzi pesanti a basso valore percepito dall’altro (forse C81 o B60) in cambio di mantenere i miei top efficienti (A08,A16,B09,B94,B67,C15,A66,A84,B75).'}\\ \tt \\ \tt PROPOSTA: {'A08', 'A16', 'B09', 'B94', 'B67', 'C15', 'A66', 'A89', 'B46', 'B12', 'B10', 'B00', 'A30', 'B60', 'C12', 'A58', 'A05', 'B85', 'C24', 'C81', 'A84', 'B75'}\\ \tt \\ \tt ARGOMENTAZIONE: {'Questa selezione resta ben entro il limite di sforzo (7816 ≤ 7868) e combina molti oggetti a basso sforzo con alta utilità. Ho incluso varietà tra A/B/C per bilanciare il rischio e lasciare spazio a scambi. Se alcuni elementi più pesanti non fossero fondamentali per te, possiamo rivedere C81 o B60, purché manteniamo la maggior parte degli elementi compatti ed efficienti come A08, A16, B09, B94, B67, C15, A66, A84 e B75.'} 
	  } 
	   } 
	   } 
	 & & \\ 
 

    \theutterance \stepcounter{utterance}  

    & & \multicolumn{4}{p{0.6\linewidth}}{\cellcolor[rgb]{0.9,0.9,0.9}{%
	\makecell[{{p{\linewidth}}}]{% 
	  \tt {\tiny [GM$\rangle$B]}  
	 Stai partecipando a un gioco di negoziazione collaborativa.\\ \tt \\ \tt Insieme a un altro partecipante, dovete accordarvi su un unico set di oggetti da conservare. Ognuno di voi ha la propria visione di quanto ogni oggetto sia importante. Non sai come l'altro partecipante valuta gli oggetti. Inoltre, ti viene fornito lo sforzo che ogni oggetto richiede.\\ \tt Potete accordarvi su un set solo se lo sforzo totale degli oggetti selezionati non supera un limite condiviso:\\ \tt \\ \tt LIMITE = 7868\\ \tt \\ \tt Ecco i valori di sforzo dei singoli oggetti:\\ \tt \\ \tt Sforzo oggetto = {"B21": 552, "B00": 560, "A08": 76, "A27": 950, "B94": 248, "B53": 871, "A23": 391, "C65": 143, "A89": 293, "B67": 208, "B45": 678, "C34": 737, "C12": 407, "A05": 366, "B60": 768, "C92": 859, "B69": 183, "A84": 231, "C06": 305, "B12": 727, "A16": 148, "B85": 357, "C35": 504, "B10": 549, "C15": 299, "B09": 91, "B16": 527, "C81": 848, "A58": 307, "C24": 214, "A30": 723, "A66": 475, "B75": 23, "B46": 298, "C70": 820}\\ \tt \\ \tt Ecco la tua visione personale sull'importanza di ogni oggetto:\\ \tt \\ \tt Valori di importanza dell'oggetto = {"B21": 138, "B00": 583, "A08": 868, "A27": 822, "B94": 783, "B53": 65, "A23": 262, "C65": 121, "A89": 508, "B67": 780, "B45": 461, "C34": 484, "C12": 668, "A05": 389, "B60": 808, "C92": 215, "B69": 97, "A84": 500, "C06": 30, "B12": 915, "A16": 856, "B85": 400, "C35": 444, "B10": 623, "C15": 781, "B09": 786, "B16": 3, "C81": 713, "A58": 457, "C24": 273, "A30": 739, "A66": 822, "B75": 235, "B46": 606, "C70": 105}\\ \tt \\ \tt Obiettivo:\\ \tt \\ \tt Il tuo obiettivo è negoziare un set condiviso di oggetti che ti avvantaggi il più possibile (cioè, massimizzi l'importanza totale per TE), rimanendo entro il LIMITEE. Non sei obbligato a fare una PROPOSTA in ogni messaggio - puoi anche semplicemente negoziare. Tutte le tattiche sono permesse!\\ \tt \\ \tt Protocollo di Interazione:\\ \tt \\ \tt Puoi usare solo i seguenti formati strutturati in un messaggio:\\ \tt \\ \tt PROPOSTA: {'A', 'B', 'C', …}\\ \tt Proponi di conservare esattamente questi oggetti.\\ \tt RIFIUTO: {'A', 'B', 'C', …}\\ \tt Rifiuta esplicitamente la proposta dell'avversario.\\ \tt ARGOMENTAZIONE: {'...'}\\ \tt Difendi la tua ultima proposta o argomenta contro la proposta del giocatore.\\ \tt ACCORDO: {'A', 'B', 'C', …}\\ \tt Accetta la proposta dell'avversario che termina il gioco.\\ \tt RAGIONAMENTO STRATEGICO: {'...'}\\ \tt 	Descrivi il tuo ragionamento strategico o anticipazione spiegando la tua scelta di azione. Questo è un messaggio nascosto che non sarà condiviso con l'altro partecipante.\\ \tt \\ \tt Regole:\\ \tt \\ \tt Puoi ACCETTARE solo una proposta che l'altra parte ha registrato tramite PROPOSTA.\\ \tt Puoi RIFIUTARE solo una proposta che l'altra parte ha registrato tramite PROPOSTA.\\ \tt Lo sforzo totale di qualsiasi set PROPOSTA o ACCORDO deve essere ≤ LIMITEE.\\ \tt NON rivelare i tuoi punteggi di importanza nascosti.\\ \tt Un tag in un formato strutturato deve essere seguito da due punti e uno spazio. L'argomento deve essere un set python contenente 0 o più stringhe.\\ \tt Quindi, deve essere della forma TAG: {...}\\ \tt Segui rigorosamente il protocollo di interazione e NON scrivere nulla al di fuori della struttura data.\\ \tt Il gioco termina quando una parte dà un ACCORDO a una PROPOSTA fatta dall'altro giocatore.\\ \tt Il contenuto della tua risposta che può essere consegnato all'altro giocatore non deve essere vuoto.\\ \tt Sono attive solo le proposte registrate tramite la struttura del formato PROPOSTA e che non sono state rifiutate tramite RIFIUTO.\\ \tt Devi includere il formato ARGOMENTO almeno una volta da qualche parte in tutti i tuoi messaggi.\\ \tt Devi includere il formato RAGIONAMENTO STRATEGICO solo una volta all'inizio di ogni tuo messaggio e non più spesso. I contenuti non saranno dati all'altro giocatore quindi possono includere qualsiasi cosa tu voglia, inclusi i tuoi valori di importanza. Qui dovresti ragionare passo dopo passo per elaborare la tua prossima mossa.\\ \tt \\ \tt Ora riceverai il primo messaggio dell'altro giocatore.\\ \tt \\ \tt PROPOSTA: {'A08', 'A16', 'B09', 'B94', 'B67', 'C15', 'A66', 'A89', 'B46', 'B12', 'B10', 'B00', 'A30', 'B60', 'C12', 'A58', 'A05', 'B85', 'C24', 'C81', 'A84', 'B75'}\\ \tt \\ \tt ARGOMENTAZIONE: {'Questa selezione resta ben entro il limite di sforzo (7816 ≤ 7868) e combina molti oggetti a basso sforzo con alta utilità. Ho incluso varietà tra A/B/C per bilanciare il rischio e lasciare spazio a scambi. Se alcuni elementi più pesanti non fossero fondamentali per te, possiamo rivedere C81 o B60, purché manteniamo la maggior parte degli elementi compatti ed efficienti come A08, A16, B09, B94, B67, C15, A66, A84 e B75.'} 
	  } 
	   } 
	   } 
	 & \\ 
 

    \theutterance \stepcounter{utterance}  

    & & & \multicolumn{4}{p{0.6\linewidth}}{\cellcolor[rgb]{1,0.85,0.72}{%
	\makecell[{{p{\linewidth}}}]{% 
	  \tt {\tiny [GM$\langle$B]}  
	 RAGIONAMENTO STRATEGICO: {'Calcolo il valore totale per me della loro proposta e verifico margine di miglioramento sostituendo elementi meno preziosi per me con altri ad alto rapporto valore/sforzo. Somma sforzi proposta: 7816 (coerente). Calcolo il mio valore: A08(868)+A16(856)+B09(786)+B94(783)+B67(780)+C15(781)+A66(822)+A89(508)+B46(606)+B12(915)+B10(623)+B00(583)+A30(739)+B60(808)+C12(668)+A58(457)+A05(389)+B85(400)+C24(273)+C81(713)+A84(500)+B75(235) = 142, (compute) Sommo: 868+856=1724; +786=2510; +783=3293; +780=4073; +781=4854; +822=5676; +508=6184; +606=6790; +915=7705; +623=8328; +583=8911; +739=9650; +808=10458; +668=11126; +457=11583; +389=11972; +400=12372; +273=12645; +713=13358; +500=13858; +235=14093. Valore 14093. Ora cerco miglioramenti: oggetti rimasti con alto valore/sforzo: B21(138/552 low), C34(484/737 ok), C35(444/504 good), B45(461/678), C70(105/820 low), C92(215/859 low), B53(65/871 bad), B16(3/527 bad), B69(97/183 low), C06(30/305 bad), C65(121/143 ok), A23(262/391 meh), A27(822/950 high value but heavy), A05 già incluso, A66 incluso, A30 incluso, C35 è molto buono per me 444/504, C65 121/143 piccolo, B21 no, B45 medio, C34 medio. Per fare spazio dobbiamo togliere qualcosa. Identifico i miei più bassi valori nella proposta: B75(235/23 è ottimo ratio comunque), C24(273/214 mediocre ma ok), A05(389/366 ok), B85(400/357 ok), A58(457/307 buono), A84(500/231 ottimo), A89(508/293 buono). I peggiori per valore per me nella lista sono forse C24(273), B75 ha valore 235 ma effort 23 quindi super efficiente, non toccarlo. Elementi non inclusi che vorrei: C35(444/504), C34(484/737), C65(121/143 piccolo), A27(822/950 grande valore), B21(138/552 no), B45(461/678 forse), B53(65 no), B16(3 no), B69(97/183 no), C06(30/305 no), C70(105/820 no), C92(215/859 no), A23(262/391 forse), A16 già, A08 già, A30 già, A66 già. Migliori aggiunte: C35 sicuramente. Forse A27 se riusciamo a togliere C81(713/848 value/effort 0.84) o elementi meno utili a me ma con alto effort: B60(808/768 very good keep), C81(713/848 ok), C12(668/407 great keep), B12(915/727 amazing keep), B10(623/549 good), B00(583/560 ok), B46(606/298 great), B94(783/248 insane), B67(780/208 insane), C15(781/299 great), A66(822/475 great), A30(739/723 borderline), A89(508/293 good), A58, A05, B85, C24, A84, B75. Candidate to drop for me: A30 (value 739, effort 723; ratio ~1.02), C81 (713/848 <1), B00 (583/560 ~1.04). If I drop C81 (848 effort), I can add C35 (504) and still have 344 effort slack, possibly also add C65 (143) and maybe A23 (391) would exceed; 504+143=647 <848, leaving +201 effort; could also upgrade A27(950) instead of C81 requires freeing extra 102 from somewhere else. Current total effort 7816. Replace C81(848) with C35(504) reduces to 7472. Then we can add A27(950) would be 8422 > limit; not possible unless we drop more. Perhaps replace A30(723) with A27(950) adds 227; need to free 227 elsewhere, maybe drop C24(214) nearly enough plus 13; could also replace B75(23) small. Better: propose replacing C81 and C24 and A30 with A27 and C35 and C65. Compute: Start 7816 - C81(848) - C24(214) - A30(723) = 7816 - 1785 = 6031. Add A27(950) + C35(504) + C65(143) = +1597 => 7628. Under limit. Value change: Remove C81(713), C24(273), A30(739) total removed value 1725. Add A27(822), C35(444), C65(121) total 1387, actually worse by 338. So not good.\\ \tt \\ \tt Alternative: Just swap C81 -> C35 and add C65 while keeping rest. New effort: 7816 - 848 + 504 + 143 = 7615. Value: -713 + 444 + 121 = -148 net, slightly worse. Not good.\\ \tt \\ \tt What if swap A30(739/723) for A27(822/950) increases value by 83 but adds 227 effort, we have only 52 slack (7868-7816=52). So need to free 175 more, maybe by swapping C24(214) for nothing or for C65(143) freeing 71, still need 104; swap B75(23) for nothing frees 23; total 94; still short 81; maybe replace A05(366) with A23(391) increases effort +25 not good. We need to drop something like B45 not in set. Maybe drop A58(307) frees 307; but A58 has value 457; net value change: +83 -457 - maybe add C65 +121 to offset. New change: +83 -457 +121 = -253 and effort change: +227 -307 +143 = +63, oops exceeds. If instead: swap A30->A27 (+227), drop C24(-214) net +13, then need -13 to fit: drop B75(-23) results net -10 fit under. Value change: +83 -273 -235 = -425, bad.\\ \tt \\ \tt Maybe the opponent's proposal is already very strong for my utility; value 14093 seems high. Perhaps accept but ask a small tweak improving my value with minimal disruption: replace C24 (low value 273) with C35 (444) if effort allows: C35 effort 504 vs C24 214 adds +290 effort; current slack 52, so would exceed by 238 unless we remove something else of at least 238 effort. We could swap B75(23) and A08(76) not good. Maybe replace B00(560, value 583) with B45(678, value 461) worse. Replace B00 with C34(737, value 484) worse for me. Replace B00 with C35(504) reduces effort by 56 and increases value -79 (583->444) actually decreases value. Not good.\\ \tt \\ \tt What about adding C65(143, value 121) replacing B75(23) increases effort +120 but we have 52 slack so would exceed. Swap C24(214) for C65(143) reduces effort by 71 and value -152; that would create extra slack to maybe upgrade A30 to A27 needing +227; with -71 slack becomes 52+71=123, still short 104.\\ \tt \\ \tt Maybe replace A30(723,739) with A66 already in; okay. Replace B00(560,583) with A27(950,822) adds +390 effort over budget. Need free 338 more (since 52 slack). Could also drop C24(214) and B75(23) total 237, still short 101; drop A58(307) then total 544 slack enough; net value change: +239 (822-583) -273 -235 -457 = -726 bad.\\ \tt \\ \tt Given this, best course is accept their proposal as it is, but maybe negotiate minor option: suggest we can lock it in. I should accept to end game maximizing.}\\ \tt \\ \tt ARGOMENTAZIONE: {'La tua proposta è molto efficiente per me e resta comodamente sotto il limite. Apprezzo in particolare il mix di elementi compatti e ad alto rendimento come A08, A16, B09, B94, B67, C15 e A66. Se per te va bene, possiamo chiudere così senza ulteriori modifiche.'}\\ \tt \\ \tt ACCORDO: {'A08', 'A16', 'B09', 'B94', 'B67', 'C15', 'A66', 'A89', 'B46', 'B12', 'B10', 'B00', 'A30', 'B60', 'C12', 'A58', 'A05', 'B85', 'C24', 'C81', 'A84', 'B75'} 
	  } 
	   } 
	   } 
	  \\ 
 

    \theutterance \stepcounter{utterance}  

    & & & \multicolumn{2}{p{0.3\linewidth}}{\cellcolor[rgb]{0.95,0.95,0.95}{%
	\makecell[{{p{\linewidth}}}]{% 
	  \tt {\tiny [GM$|$GM]}  
	 game successful 
	  } 
	   } 
	   } 
	 & & \\ 
 

    \theutterance \stepcounter{utterance}  

    & & & \multicolumn{2}{p{0.3\linewidth}}{\cellcolor[rgb]{0.95,0.95,0.95}{%
	\makecell[{{p{\linewidth}}}]{% 
	  \tt {\tiny [GM$|$GM]}  
	 end game 
	  } 
	   } 
	   } 
	 & & \\ 
 

\end{supertabular}
}

\end{document}
